%% ma: L12-B02-0015
%% lop: 12
%% chuong: 1
%% bai: 2
%% dang: 12.2.4
%% loai: TN4
%% muc: TH
%% dap_an: "D"
%% nguon: "Ảnh giáo viên gửi 05/11/2026 (M11): Kiểm tra giữa HK1 Toán 12 (2026-2027), THPT Hồ Thị Bi (TP.HCM), mã đề 1133, Phần I Câu 3"
%% kien_thuc: "Bài 2 (giá trị nhỏ nhất của hàm số trên đoạn đọc từ bảng biến thiên)"
%% trang_thai: chinh-thuc
%% ghi_chu: "Đề gốc không kèm đáp án; đã giải độc lập. Dạng 12.2.4 gốc hỏi tổng M + m; câu này hỏi giá trị nhỏ nhất, cùng mô hình đọc từ bảng biến thiên. kho.py trung gợi ý thay số so với L12-B02-0007 (khung 0,80) nhưng khác bảng biến thiên và hỏi giá trị nhỏ nhất; giữ lại."
\begin{ex}
Cho hàm số $y=f(x)$ liên tục trên $\mathbb R$ và có bảng biến thiên như dưới đây:
\begin{center}
\begin{tikzpicture}
  \tkzTabInit[lgt=1.5,espcl=2.4]{$x$/.8,$y'$/.8,$y$/2.4}{$-\infty$,$-1$,$1$,$+\infty$}
  \tkzTabLine{,+,0,-,d,+,}
  \tkzTabVar{-/$-2$,+/$5$,-/$0$,+/$1$}
\end{tikzpicture}
\end{center}
Giá trị nhỏ nhất của hàm số $y=f(x)$ trên đoạn $[-1;1]$ bằng
\choice{$-1$}{$5$}{$-2$}{\True $0$}
\loigiai{Trên đoạn $[-1;1]$, $f'(x)<0$ với $-1<x<1$ nên hàm số nghịch biến.
Do đó $\min\limits_{[-1;1]}f(x)=f(1)=0$.}
\end{ex}
